\section{Discussion}
\label{sec:discussion}

\paragraph{Why search wins here.} The learned policy builds a plan in one pass of sequential decisions, and sampling
more rollouts adds diversity but no improvement of a given plan. ALNS spends the same time improving one plan with
many cheap, exact moves: our representation keeps every move feasible, and Eq.~\eqref{eq:insert} scores an insertion
in time linear in the coalition size, so a single core evaluates thousands of neighbours per second. The CP-SAT
baselines search a richer neighbourhood per step but pay for building and solving a model at every step, which is why
parallelising them helps little.

\paragraph{Implications.} We do not claim that learning is the wrong tool for multi-robot coordination: a
decentralised policy can react to events that a central plan cannot. But the static planning performance of
learned policies should be measured against strong search at matched compute, on the same cores and with the CPU
use reported. The published comparison missed plans \ConeBrlgainlo--\ConeBrlgainhi\% shorter, which a standard
operations-research technique adapted to coalitions finds at the same budget, and largely already on one core in 2\,s.

\paragraph{Limitations.}
\begin{itemize}
\item \emph{Static, centralised planning.} We evaluate plans for fully known, static instances, as the benchmark
  does. The RL policy is decentralised and reactive, which matters under uncertainty and with dynamic arrivals; we
  did not evaluate such settings in the confirmatory run. An exploratory dynamic variant did not pass the go/no-go
  gates we had set on dev, and it is reported in the supplement as a negative result.
\item \emph{One benchmark.} All instances come from one generator with Euclidean travel, one speed and a makespan
  objective. We expect the representation to transfer to other objectives and travel models, but have not tested it.
\item \emph{RL on CPU.} To match compute, the policy runs on the same CPU cores as every other method. On a GPU it
  could draw more rollouts. It already draws \RLNlo{} to \RLNhi{} times the published number, which makes its plans
  only \RLsamplegainlo--\RLsamplegainhi\% shorter than with ten rollouts.
\item \emph{Unknown optimality gaps.} On 50 or more tasks no method proves optimality, and the certified bounds are
  loose, so we report gaps to the best known plans and ratios between methods.
\item \emph{Tuning.} ALNS was developed and tuned on the dev split of the same generator. Every competitor with a
  parameter was tuned per setting on the same split, which favours the competitors.
\item \emph{Pre-registration.} The val dry run was known before the freeze, the test-split results of the
  constructor and of the RL policy had been seen in an earlier phase, and the released test set of the benchmark
  (not blind) was used to choose CTAS-D's solver. No ALNS or CP-SAT variant and no CTAS-D port ran on the test split
  before the freeze.
\item \emph{Single host and seed.} All timed runs used one host and one seed per instance, method and budget. On
  val, 52 cells measured on two different hosts agreed within 2\%.
\item \emph{CTAS-D.} Our port solves the published model with an open solver on a time grid of $10^{-3}$ and counts
  model building in the budget; it is not the benchmark paper's Gurobi run, although it beats that run on the
  released 20-task instances.
\end{itemize}

\section{Conclusion}
\label{sec:conclusion}

On the HeteroMRTA benchmark, an adaptive large neighbourhood search over a deadlock-free coalition representation
with exact linear-time insertion finds plans \ConeBrlgainlo--\ConeBrlgainhi\% shorter than the published
reinforcement-learning policy and \ConeBclasslo--\ConeBclasshi\% shorter than three CP-SAT baselines and restarts, while
the exact MILP rarely finds a plan at all, at equal compute and under a pre-registered protocol. We release the solver, the harness, the pre-registration and every result row, so that future learned
planners can be compared against this baseline at matched budgets.
